% Generated draft: requires mathematical and source review.
\documentclass[11pt]{article}
\usepackage[margin=25mm]{geometry}
\usepackage{fontspec}
\setmainfont{Latin Modern Roman}
\usepackage{amsmath,amssymb,mathtools,unicode-math}
\setmathfont{Latin Modern Math}
\usepackage{enumitem,needspace,xurl,hyperref}
\setlength{\parindent}{0pt}
\setlength{\parskip}{5pt}
\setlength{\emergencystretch}{3em}
\newcommand{\N}{\mathbb{N}}
\newcommand{\Z}{\mathbb{Z}}
\newcommand{\Q}{\mathbb{Q}}
\newcommand{\R}{\mathbb{R}}
\newcommand{\C}{\mathbb{C}}
\newcommand{\F}{\mathbb{F}}
\newcommand{\A}{\mathbb{A}}
\newcommand{\PP}{\mathbb{P}}
\newcommand{\E}{\mathbb{E}}
\newcommand{\eps}{\varepsilon}
\newcommand{\dd}{\,\mathrm{d}}
\begin{document}
\section*{UnsolvedMath ID: 20003149}
\textbf{Source problem code:} catalogue entry
\subsection{Definitions and mathematical statement}
For a principal \(G\)-bundle \(P\to B\) with connection curvature \(\Omega\in\Omega^2(B;\operatorname{Ad}P)\), use the standard Weinstein meaning of fatness: for every nonzero vertical Lie-algebra covector, the paired curvature is nondegenerate on the horizontal tangent space. With an invariant inner product this can be written using nonzero \(u\in\mathfrak g\). The AIM question asks whether existence of a fat principal bundle forces \(G\cong S^1, S^3\) or \(SO(3)\). It does not explicitly require positive or nonnegative curvature of an entire connection-metric total space.

The supplied aid reports compact connected torus counterexamples under that literal definition. Let \(B=T^4=(\R/\Z)^4\), \(\mathrm{vol}=\dd x_1\wedge\dd x_2\wedge\dd x_3\wedge\dd x_4\), and put
\[
\omega_1=\dd x_1\wedge\dd x_2+\dd x_3\wedge\dd x_4,\quad
\omega_2=\dd x_1\wedge\dd x_3-\dd x_2\wedge\dd x_4,\quad
\omega_3=\dd x_1\wedge\dd x_4+\dd x_2\wedge\dd x_3.
\]
These are closed integral forms satisfying
\[
\omega_i\wedge\omega_j=2\delta_{ij}\,\mathrm{vol},\qquad
\left(\sum_{i=1}^ru_i\omega_i\right)^2=2\sum_{i=1}^ru_i^2\,\mathrm{vol}\ne0
\quad(u\ne0),
\]
for \(r=2,3\). Prequantize the first \(r\) classes by principal circles \(\R/(2\pi\Z)\) with real connection curvature \(2\pi\omega_i\), and take their fiber product. The resulting principal \(T^r\)-bundle has curvature \(2\pi\sum_{i=1}^r\omega_i e_i\); every nonzero pairing is symplectic, so the connection is fat. Thus \(T^2\) and \(T^3\) disprove the stated group list. The same reasoning applies to integral rank-two subpencils; an arbitrary real subpencil still requires the period/lattice integrality needed for prequantization.

The aid also audits the flat-base connection metric: its abelian torus fibers have zero vertical sectional curvature, and a horizontal plane with nonzero curvature form has negative sectional curvature by O'Neill's formula. Hence this particular total-space metric is neither positively nor everywhere nonnegatively curved. Fatness governs the relevant mixed directions and does not imply positive curvature of every plane. If the workshop intended additional total-curvature hypotheses, these examples do not settle that stronger reading. The report refutes literal fatness alone and does not claim a classification with unspecified curvature assumptions.
\subsection{Short English statement}
Under standard fatness alone, the supplied integral symplectic pencil yields compact fat torus bundles, contradicting the three-group list. Their flat-base connection metrics have zero vertical and negative horizontal planes, so the examples do not answer a stronger question imposing total-space curvature.
\subsection{Sources}
\begin{itemize}
\item Problem 24, AIM workshop \emph{Manifolds with nonnegative sectional curvature}, \url{https://aimath.org/WWN/nnsectcurvature/nnsectcurvature.pdf}.
\item The research summary identifies its fatness convention as Weinstein's, as used by Ziller and Florit--Ziller, and invokes O'Neill's curvature formula. Full bibliographic entries are not supplied.
\item Local UnsolvedMath registry, permanent ID 20003149, exact code \texttt{catalogue entry}; canonical location \texttt{aim-topology-notes.json}, notes[236], dated 14 August 2026; statement labeled exact, registry status solved and research classification solved in literature.
\item Research attempt 1: counterexample, problem refuted at the run, rigor label proved, novelty confidence low. The normalized torus construction and curvature audit are supplied machine-generated research aids, not peer reviewed; independent verification is required before citation.
\end{itemize}
\end{document}
